\section{A depth-independent bound and a mixed-contact example}
\label{sec:mixed}
Fix the local field and uniformizer as above. Let $j$ range from $1$ to the
largest contact depth $H\ge1$, set $m_j=\#\{c\ne b:v_\pi(b-c)\ge j\}$, and write
$m=m_H$, $h=m-1$, $e=v_\pi(p)$, $D_j=(q-1)q^{j-1}$ and $D=D_H$.
Write $h_c=v_\pi(b-c)$ and $\omega=\sum_{c\ne b}q^{h_c}$. Set $n=r-1$ and $\Delta(n)=\omega+Dn-t_b(n+1)$. The period questions for
$\Delta(n)$ and $t_b(r)-Dr$ are identical.
The earlier ProveIt paper \emph{Prime patterns and branch precision}
gives a fixed-ring recurrence and a period bound containing $q^{H-1}$.
\href{https://github.com/VladimirReshetnikov/ProveIt/blob/b9a7c196e4aed3d93d4f222efef677a50efe912b/Analysis/GeneralizedRiemannHypothesis/Research/grh_prime_patterns_and_branch_precision/periodicity.tex}{See its section on the exact defect.}
We give a smaller universal test window whose length and starting index are independent of $H$ and of the counts $m_j$ for $j<H$.

\begin{theorem}
For $m=1$ the defect is identically zero. For $m\ge2$ put
\[
 L=p^{\lfloor\log_p h\rfloor+\lceil h/e\rceil},
 \qquad n_0=\left\lceil\frac{qh}{q-1}\right\rceil-1.
\]
Then $0\le\Delta(n)\le Dh$ at all precisions and
$\Delta(n+L)=\Delta(n)$ for all $n\ge n_0$.
In particular, the least eventual period is the least power of $p$ whose
cyclic shift fixes the $L$ entries starting at $n_0$.
\end{theorem}

\begin{proof}
We give the finite coefficient test explicitly.
Let $f(X)=\pi X+X^q$, $g_j(X)=\pi+(f^{\circ(j-1)}(X))^{q-1}$ and
let $d$ denote the total number of directions, including $b$.
Up to an integral unit the denominator is
$P=X^{d-1}\prod_jg_j^{m_j}$, of degree $\omega$.
Define
\[
 A_j(X)=X^{D-D_j}(g_j(X)-X^{D_j})/\pi.
\]
These are integral polynomials. The constant coefficient of $A_H$ is $1$,
and $A_j$ is divisible by $X^{D-D_j}$.
We have $g_j=X^{D_j}(1+\pi X^{-D}A_j)$.
In precision $k+1$ the coefficient of $X^{-\omega-Dk}$ in $1/P$
is $(-\pi)^k\binom{k+h}{h}$, since the smallest exponent requires use of the
deepest factor alone. No smaller exponent survives. Thus a binomial
which is a unit gives $t_b(k+1)=\omega+Dk$.
Among any $h+1$ consecutive values $\binom{k+h}{h}$ for $k\ge0$ there is
a unit, since the $h$th forward difference is $1$. The value at $k=0$
also works. Reduction from a higher precision to a lower one gives monotonicity
of $t_b$. Choosing such a $k\in[\max(0,n-h),n]$ proves the defect bound,
including the case $h=0$.

For $h>0$, multiply $1/P$ by $(-1)^n X^{\omega+Dn}A_H^{-n}$ in
$(\mathcal O/\pi^{n+1})((X))$. The result is a power series of order
$\Delta(n)$, since $A_H$ is a unit. Removing this unit power is the step
which eliminates the depth from the period bound.
The negative-binomial expansion of the factors has an index $k$ at $H$
and indices $s_j\ge0$ at $j<H$. Put $u=\sum_{j<H}s_j$.
All terms with $k+u>n$ vanish. Set $\ell=n-k-u$ and $v=\ell+u$.
The corresponding term in the power series is
\[
 (-1)^\ell\pi^{n-\ell}\binom{n-v+h}{h}X^{D\ell}A_H^{-v}
          \prod_{j<H}\binom{m_j+s_j-1}{s_j} A_j^{s_j}.
\]
Its order is at least $D\ell+\sum_{j<H}(D-D_j)s_j$.
As $D_j\le D/q$ for $j<H$, a contribution to degree $a<Dh$
requires $v<qh/(q-1)$, hence $v\le n_0$. For $n\ge n_0$ all
these candidate tuples have $k=n-v\ge0$.

Fix such a degree $a$ and let $b_a=\lfloor a/D\rfloor\le h-1$.
All contributing terms have $\ell\le b_a$. Divide the coefficient
at $a$ in the displayed expansion by $\pi^{n-b_a}$ and test it
modulo $\pi^{b_a+1}$. This is exactly its survival test modulo $\pi^{n+1}$.
After division every summand depends on $n$ only through
$\binom{n-v+h}{h}$.
For $1\le i\le h$ we have
\[
 v_p\binom{L}{i}\ \ge\ v_p(L)-v_p(i)\ \ge\ \lceil h/e\rceil,
\]
by $i\binom{L}{i}=L\binom{L-1}{i-1}$. Vandermonde's identity now
makes $\binom{n-v+h}{h}$ periodic modulo $\pi^h$ with period $L$.
This proves periodicity of every tested coefficient's survival. If all
degrees below $Dh$ vanish, the defect bound forces $\Delta(n)=Dh$.
The last assertion follows because the stabilizer of a cyclic word of
length $L$ is a subgroup of $\mathbb Z/L\mathbb Z$.
\end{proof}

When $h=m-1$ is a power of $p$, including $h=1$, the period bound is sharp at every ramification index: the depth-one formula \eqref{eq:period} attains the exponent $\lfloor\log_p h\rfloor+\lceil h/e\rceil$ at $j=h$.

\paragraph{The shallow branches can change the least period.}
Over $\mathbb Q_2$, take five other directions at depth two. Compare this
with the same configuration augmented by two directions at depth one.
For example use $b=0$, take $4,12,20,28,36$, and then also $2,6$.
Both profiles have $H=2$, $m=5$, $D=2$, $n_0=7$ and $L=64$.
Their least eventual periods are respectively $32$ and $64$.
Here is a finite calculation certifying the claim. With $j$ the least
nonnegative residue of $n$ modulo $16$, the entries below give the
defect periodic in $n$ from $n=7$. In either row the entries for
$j=0,1,2,3,8,9,10,11$ are zero and those for $j=4,12$ are one.
\[
\begin{array}{c|rrrrrr}
 & j=5&j=6&j=7&j=13&j=14&j=15\\\hline
 (2,2,2,2,2)&3&2&2&2&3&4,5\\
 (2,2,2,2,2,1,1)&2&2&2&3&4,5&4,6,4,7
\end{array}
\]
A list in a cell is read in order as $\lfloor n/16\rfloor$ increases from
zero, cyclically. For instance the second row takes different values at
$n=31$ and $63$, which rules out period $32$. To reproduce each entry, use
$P=\prod_{c\ne b} f^{\circ h_c}(X)$ and reduce $1,X,X^2,\ldots$ modulo
$(P,2^{n+1})$ until the first zero. If its exponent is $t$, record
$\deg P+2n-t$. Checking the 64 integers $7\le n\le70$ suffices by the
theorem. The accompanying program \texttt{checks/mixed\_period.py} performs
this exact computation and cross-checks it by the scaled reciprocal recurrence.
